\begin{frame}[t]{Da linguagem C ao ambiente Arduino}
\begin{itemize}
\item Nas aulas anteriores: programas C11 no computador, com \texttt{main} e \texttt{printf}.
\item Nesta prática: arquivos \texttt{.ino} compilados como C++ pelo ambiente Arduino.
\item Variáveis, operadores, condicionais, laços e funções continuam presentes.
\item O ambiente fornece a inicialização e chama \texttt{setup} e \texttt{loop}.
\item Funções de GPIO fazem a ponte entre a lógica e a placa.
\end{itemize}
\begin{block}{Diferença importante}
\texttt{bool} faz parte do C++; nos exemplos C11 usamos \texttt{stdbool.h}. O código Arduino não é simplesmente um programa C de terminal.
\end{block}
\end{frame}
